\chapter[APRESENTAÇÃO DA TRAJETÓRIA DO ALUNO]{APRESENTAÇÃO DA TRAJETÓRIA DO ALUNO}

Minha trajetória na área de \ac{ti} teve início em agosto de 2021, quando ingressei no novo curso da \ac{pucrs} de Ciência de Dados e Inteligência Artificial. Anteriormente, havia concluído um curso de Bacharelado em Química Industrial na mesma universidade, mas percebi que minha afinidade com essa área já não era mais tão forte quanto na época em que era calouro. Decidi, então, embarcar em uma nova jornada acadêmica, optando pelo campo que mais despertava minha curiosidade e meu interesse.

Após completar três semestres e perceber que a pesquisa acadêmica e
disciplinas teóricas não me cativavam mais, a Engenharia de Software me chamou atenção. Essa mudança de foco me levou a trocar de curso no final de 2022. Embora mantenha um afeto pela área de dados e pelas tecnologias que estudei nos primeiros meses do curso, foi o desejo de participar ativamente na concepção, construção e manutenção de software que impulsionou essa decisão.

No mesmo semestre em que troquei de curso, consegui me inserir no mercado de trabalho como estagiário na empresa PORTOSYS, terceirizada da Industrial and Financial Systems.

Atualmente, sigo nesse estágio. Nele, desenvolvo soluções de \ac{erp}. A principal tecnologia utilizada nesse trabalho é \ac{plsql}. Também emprego C\# e \ac{sql}.

Este memorial apresenta, nos capítulos seguintes, minha trajetória nas quatro etapas da \ac{ages}.
